\titre{}
\theme{limites}
\auteur{Nathan Scheinmann}
\niveau{3M}
\source{sesamath3e}
\type{serie}
\piments{2}
\pts{}
\annee{2425}

\contenu{
\tcblower
\noindent Calculer les limites suivantes en justifiant chaque étape.  
\begin{tasks}(3)
\task \(\displaystyle\lim_{x\to -2}\frac{-2}{2x+4}\)
\task \(\displaystyle\lim_{x\to 1}\frac{x-2}{x^2+3x-4}\)
\task \(\displaystyle\lim_{x\to 2}\frac{3x-7}{x-2}\)
\task \(\displaystyle\lim_{x\to -2}\frac{-3x+2}{x+2}\)
\task \(\displaystyle\lim_{x\to 1}\frac{3x}{1-x^2}\)
\task \(\displaystyle\lim_{x\to -3}\frac{5x}{x+3}\)
\end{tasks}
}
\correction{
	\tcblower

\begin{tasks}

  \task \( \displaystyle \lim_{x \to -2} \dfrac{-2}{2x + 4}
  \stackrel{\text{P4}}{=} \dfrac{-2}{\lim_{x \to -2} (2x + 4)}
  \stackrel{\text{P1, L2}}{=} \dfrac{-2}{2\cdot(-2) + 4}
  = \dfrac{-2}{-4 + 4} = \dfrac{-2}{0} \)
  
  Type \( \dfrac{1}{0} \) : limite infinie → on doit préciser le signe du dénominateur :
  
  \[
  \lim_{x \to -2^-} \dfrac{-2}{2x + 4} = \dfrac{-2}{0^-} = +\infty, \quad
  \lim_{x \to -2^+} \dfrac{-2}{0^+} = -\infty
  \]
  La limite n’existe pas.


  \task \( \displaystyle \lim_{x \to 1} \dfrac{x - 2}{x^2 + 3x - 4}
  = \dfrac{1 - 2}{1^2 + 3\cdot1 - 4}
  = \dfrac{-1}{0} \)

  Type \( \dfrac{1}{0} \) → on factorise : \( x^2 + 3x - 4 = (x + 4)(x - 1) \)  
  \[
  \lim_{x \to 1^-} \dfrac{x - 2}{(x + 4)(x - 1)} = \dfrac{-1}{(5)(0^-)} = -\infty, \quad
  \lim_{x \to 1^+} = \dfrac{-1}{(5)(0^+)} = -\infty
  \]
  Donc \( \lim_{x \to 1} f(x) = -\infty \)

  \task \( \displaystyle \lim_{x \to 2} \dfrac{3x - 7}{x - 2}
  = \dfrac{6 - 7}{0} = \dfrac{-1}{0} \)

  Limite infinie :  
  \[
  \lim_{x \to 2^-} = \dfrac{-1}{0^-} = +\infty, \quad
  \lim_{x \to 2^+} = \dfrac{-1}{0^+} = -\infty
  \]
  Donc la limite n’existe pas.

  \task \( \displaystyle \lim_{x \to -2} \dfrac{-3x + 2}{x + 2}
  = \dfrac{6 + 2}{0} = \dfrac{8}{0} \)

  Limite infinie :
  \[
  \lim_{x \to -2^-} = \dfrac{8}{0^-} = -\infty, \quad
  \lim_{x \to -2^+} = \dfrac{8}{0^+} = +\infty
  \]
  Donc la limite n’existe pas.

  \task \( \displaystyle \lim_{x \to 1} \dfrac{3x}{1 - x^2}
  = \dfrac{3}{1 - 1} = \dfrac{3}{0} \)

  \( 1 - x^2 = -(x - 1)(x + 1) \), donc :
  \[
  \lim_{x \to 1^-} \dfrac{3x}{1 - x^2} = \dfrac{3}{-(0^-)(2)} = -\infty, \quad
  \lim_{x \to 1^+} = \dfrac{3}{-(0^+)(2)} = -\infty
  \]
  Donc la limite est \( -\infty \).

  \task \( \displaystyle \lim_{x \to -3} \dfrac{5x}{x + 3} = \dfrac{-15}{0} \)

  Limite infinie :
  \[
  \lim_{x \to -3^-} = \dfrac{-15}{0^-} = +\infty, \quad
  \lim_{x \to -3^+} = \dfrac{-15}{0^+} = -\infty
  \]
  Donc la limite n’existe pas.

\end{tasks}
}

